\paragraph{带符号圆维：两侧单位与单侧正锥的分账}\label{para:cdim-signed-circle-dimension}
定义 \ref{def:half-circle-dimension} 的半圆维 \(\cdim_{+}(M)\) 以群完成的圆维计量单侧锥的容量。
当 \(M\) 同时含有两侧可逆的单位元方向时，半圆维将该方向也以“半”计入；为把两侧与单侧通道在同一计量中结构性分账，引入如下半整数不变量。

\begin{definition}[带符号圆维]\label{def:cdim-signed}
令 \(M\) 为有限生成、可消去、交换幺半群，并满足定义 \ref{def:half-circle-dimension} 的实现条件。
记其单位元群为 \(U(M)\subseteq M\)，并将其视为群完成 \(G(M)\) 的子群。
定义 \(M\) 的\textbf{带符号圆维}为
\[
\boxed{\;
\cdim^{\pm}(M):=\cdim\!\bigl(U(M)\bigr)+\frac12\,\cdim\!\bigl(G(M)/U(M)\bigr)
\in\tfrac12\NN\cup\{0\}\; } .
\]
\end{definition}

\begin{proposition}[带符号圆维的延拓与稳定性]\label{prop:cdim-signed-laws}
在定义 \ref{def:cdim-signed} 的设定下：
\begin{enumerate}
  \item 若 \(M\) 本身为群，则 \(\cdim^{\pm}(M)=\cdim(M)\)。
  \item 若 \(U(M)\) 有限，则 \(\cdim^{\pm}(M)=\cdim_{+}(M)\)。
  \item 对任意同类幺半群 \(M,N\)，有 \(\cdim^{\pm}(M\oplus N)=\cdim^{\pm}(M)+\cdim^{\pm}(N)\)。
  \item 设 \(f:M'\to M\) 为幺半群同态。
        若诱导的群同态 \(U(M')\to U(M)\) 与 \(G(M')/U(M')\to G(M)/U(M)\) 均为有限扩张（有限核且余核有限），
        则 \(\cdim^{\pm}(M')=\cdim^{\pm}(M)\)。
\end{enumerate}
\leanverified{paper\_cdim\_signed\_laws}
\end{proposition}
\begin{proof}
(1) 当 \(M\) 为群时 \(U(M)=M\) 且 \(G(M)/U(M)=0\)，代入定义即得。
(2) 若 \(U(M)\) 有限，则 \(\cdim(U(M))=0\)（例 \ref{ex:circle-dimension-examples}），并且 \(G(M)\to G(M)/U(M)\) 为有限扩张，故 \(\cdim(G(M)/U(M))=\cdim(G(M))\)（命题 \ref{prop:circle-dimension-laws}(3)）。
代回得到 \(\cdim^{\pm}(M)=\frac12\cdim(G(M))=\cdim_{+}(M)\)。
(3) 对可消去交换幺半群有 \(G(M\oplus N)\cong G(M)\oplus G(N)\) 且 \(U(M\oplus N)\cong U(M)\oplus U(N)\)；
再注意
\(
G(M\oplus N)/U(M\oplus N)\cong (G(M)/U(M))\oplus (G(N)/U(N))
\)，
并用命题 \ref{prop:circle-dimension-laws}(2) 即得。
(4) 由有限扩张不变性（命题 \ref{prop:circle-dimension-laws}(3)）分别作用于两项并相加即可。证毕。
\end{proof}

\begin{proposition}[正交幺半群的闭式计算]\label{prop:cdim-signed-orthant-closed}
设
\[
M=\ZZ^{u}\oplus \NN^{v}\oplus F,
\]
其中 \(F\) 为有限交换群。
则
\[
\boxed{\;\cdim^{\pm}(M)=u+\frac{v}{2}\;}
\]
并且 \(\cdim_{+}(M)=\frac12(u+v)\)。
\leanverified{paper\_cdim\_signed\_orthant\_closed}
\end{proposition}
\begin{proof}
在加法幺半群意义下，\(\NN^{v}\) 的单位元仅为 \(0\)，故
\(
U(M)\cong \ZZ^{u}\oplus F
\)；
而群完成为
\(
G(M)\cong \ZZ^{u+v}\oplus F
\)，
从而
\(
G(M)/U(M)\cong \ZZ^{v}
\)。
代入定义 \ref{def:cdim-signed} 并用例 \ref{ex:circle-dimension-examples} 得到
\(
\cdim^{\pm}(M)=\cdim(\ZZ^{u})+\frac12\cdim(\ZZ^{v})=u+v/2
\)。
另一方面 \(\cdim_{+}(M)=\frac12\cdim(G(M))=\frac12(u+v)\)。证毕。
\end{proof}

\begin{proposition}[Gödel 标量化的位长必含 \texorpdfstring{$\log B$}{logB} 因子]\label{prop:cdim-godel-scalarization-bitlength-logB}
令 \(B\ge 1\)，并令加法自由交换幺半群 \(M_B:=\NN^{B}\)。
则 \(\cdim^{\pm}(M_B)=\frac{B}{2}\)。
设 \(\iota:M_B\to \NN\) 为任意注入的幺半群同态，且存在互异素数 \(p_1,\dots,p_B\) 与整数 \(a_i\ge 1\) 使
\[
\iota(v_1,\dots,v_B)=\prod_{i=1}^{B}p_i^{a_i v_i}.
\]
则在超立方体顶点集合 \(\{0,1\}^{B}\subset M_B\) 上有
\[
\max_{v\in\{0,1\}^B}\log_2\iota(v)=\Theta(B\log_2 B)
=\Theta\!\bigl(\cdim^{\pm}(M_B)\,\log \cdim^{\pm}(M_B)\bigr).
\]
因此，将 \(B\) 个边界自由度完全标量化为单个整数时，最坏情形位长在 \(B\) 的线性尺度之外必带有额外的 \(\log B\) 乘子。
\leanverified{paper\_cdim\_godel\_scalarization\_bitlength\_logB}
\end{proposition}
\begin{proof}
取 \(v=(1,\dots,1)\)，则
\(
\iota(v)=\prod_{i=1}^{B}p_i^{a_i}\ge \prod_{i=1}^{B}p_i
\)。
记第 \(i\) 个素数为 \(\mathfrak p_i\)，并记 primorial \(\mathfrak p_B\#:=\prod_{i=1}^{B}\mathfrak p_i\)。
由于任意 \(B\) 个互异素数的乘积不小于 \(\mathfrak p_B\#\)，并且 \(\log(\mathfrak p_B\#)=\Theta(B\log B)\)，得到下界
\(
\log_2\iota(v)\ge c\,B\log_2 B
\)。
从而 \(\max_{v\in\{0,1\}^B}\log_2\iota(v)\ge c\,B\log_2 B\)。
\par
反向上界：令 \(\iota_0(v)=\prod_{i=1}^{B}\mathfrak p_i^{v_i}\)，则对任意 \(v\in\{0,1\}^B\) 有
\(
\log \iota_0(v)\le \sum_{i=1}^{B}\log \mathfrak p_i=O(B\log B)
\)。
故最坏情形位长至多为 \(C B\log B\)，合并得 \(\Theta(B\log B)\)。
最后以 \(\cdim^{\pm}(M_B)=B/2\) 代入即可得到等价表述。证毕。
\end{proof}

\paragraph{Stokes 角色空间与相对上同调：带符号圆维的拓扑刻画}\label{para:cdim-signed-stokes-cohomology}
命题 \ref{prop:cdim-signed-orthant-closed} 给出了 \(\cdim^{\pm}\) 在正交幺半群上的闭式。
相应地，可将单位元方向与单侧正锥方向分别落实为环面与圆盘，并以配对上同调的连接同态表达边界--体积转移，从而在同一体系内得到与 \(\cdim^{\pm}\) 一致的拓扑刻画。

\begin{definition}[Stokes 角色空间与边界角色空间]\label{def:cdim-stokes-character-pair}
令 \(M\) 为有限生成、可消去交换幺半群。
记闭单位圆盘 \(\overline{\mathbb D}:=\{z\in\CC:\ |z|\le 1\}\) 与单位圆 \(\TT\) 均取乘法幺半群结构。
定义
\[
\mathfrak X(M):=\Hom_{\mathrm{CMon}}(M,\overline{\mathbb D}),
\qquad
\Sigma(M):=\Hom_{\mathrm{Ab}}\!\bigl(G(M),\TT\bigr)=\widehat{G(M)}.
\]
将 \(\mathfrak X(M)\) 与 \(\Sigma(M)\) 赋以点态收敛拓扑，并记
\(\mathfrak X_0(M)\subseteq \mathfrak X(M)\)、\(\Sigma_0(M)\subseteq \Sigma(M)\)
为包含平凡角色的连通分支。
\end{definition}

\begin{definition}[Stokes 加权上同调维数]\label{def:cdim-stokes-weighted-dimension}
在定义 \ref{def:cdim-stokes-character-pair} 的设定下，定义
\[
\mathrm{wdim}(M)
:=\rank_{\ZZ}H^{1}\bigl(\mathfrak X_0(M);\ZZ\bigr)
\;+\;\frac12\,\rank_{\ZZ}H^{2}\bigl(\mathfrak X_0(M),\Sigma_0(M);\ZZ\bigr)
\in\tfrac12\NN\cup\{0\}.
\]
\end{definition}

\begin{theorem}[一般可消交换幺半群的 Stokes 角色空间收缩与 \texorpdfstring{$\mathrm{wdim}$}{wdim} 的完全刻画]\label{thm:cdim-stokes-character-contraction-general-monoid}
令 \(M\) 为有限生成、可消去、交换幺半群。
记单位群 \(U:=U(M)\)、群完成 \(G:=G(M)\)，并令
\[
u:=\rank_{\ZZ}U,\qquad v:=\rank_{\ZZ}(G/U).
\]
则存在幺半群同态 \(\ell:M\to\NN\) 使得 \(\ell^{-1}(0)=U\)，并且由
\[
H:[0,1]\times \mathfrak X_0(M)\to \mathfrak X_0(M),\qquad
(t,\chi)\longmapsto \bigl(m\mapsto \chi(m)\,t^{\ell(m)}\bigr)
\]
给出的同伦 \(H\) 将 \(\mathfrak X_0(M)\) 强变形收缩到子空间
\[
\mathfrak X_U:=\{\chi\in \mathfrak X_0(M):\ \chi(m)=0\ \text{对一切 }m\notin U\}.
\]
并且自然同胚
\[
\mathfrak X_U\ \cong\ \Hom_{\mathrm{Ab}}(U,\TT)_0\ \cong\ \TT^{u},
\]
从而
\[
H^{1}(\mathfrak X_0(M);\ZZ)\cong \ZZ^{u}.
\]
此外，包含映射 \(i:\Sigma_0(M)\hookrightarrow \mathfrak X_0(M)\) 在 \(H^{1}\) 上诱导单射，
并且相对上同调群满足自然同构
\[
H^{2}\bigl(\mathfrak X_0(M),\Sigma_0(M);\ZZ\bigr)\ \cong\ \Hom(G/U,\ZZ),
\]
从而
\[
\boxed{\;\mathrm{wdim}(M)=u+\frac{v}{2}=\cdim^{\pm}(M).\;}
\]
\leanverified{paper\_cdim\_stokes\_character\_contraction\_general\_monoid}
\end{theorem}
\begin{proof}[证明要点]
令 \(P:=M/U\)。
由于 \(M\) 可消且有限生成，\(P\) 仍为有限生成可消交换幺半群，且无非平凡单位。
将 \(P\) 嵌入其群完成的自由部分 \(G(P)_{\mathrm{free}}\cong \ZZ^{v}\)，并令
\(
C:=\RR_{\ge 0}\cdot P\subset \RR^{v}
\)
为其生成的有理多面体锥。
因 \(P\) 无单位，\(C\) 为尖锥，故其对偶锥 \(C^{\vee}\) 内点非空。
取内点 \(\lambda\in C^{\vee}\cap \ZZ^{v}\)，则 \(\lambda(p)>0\) 对一切 \(p\in P\setminus\{0\}\)。
令 \(\ell_P:=\lambda|_{P}:P\to\NN\)，并令 \(\pi:M\to P\) 为商映射，取 \(\ell:=\ell_P\circ\pi\)，则 \(\ell^{-1}(0)=U\)。
\par
对任意 \(t\in[0,1]\)，映射 \(a_t(m):=t^{\ell(m)}\) 是 \(M\to\overline{\mathbb D}\) 的幺半群同态，故 \(H_t(\chi):=\chi\cdot a_t\in\mathfrak X_0(M)\)。
当 \(t=0\) 时，\(\ell(m)>0\) 等价于 \(m\notin U\)，从而 \(H_0(\chi)(m)=0\)（\(m\notin U\)）且 \(H_0(\chi)(m)=\chi(m)\)（\(m\in U\)），故像落入 \(\mathfrak X_U\)；
并且对 \(\chi\in\mathfrak X_U\) 有 \(H_t(\chi)=\chi\)，故 \(H\) 为强变形收缩。
单位群的像必落在 \(\TT\)，因此 \(\mathfrak X_U\cong \Hom(U,\TT)_0\cong \TT^{u}\)，得到 \(H^{1}(\mathfrak X_0(M);\ZZ)\cong \ZZ^{u}\)。
\par
另一方面，\(\Sigma_0(M)\cong \Hom(G,\TT)_0\cong \TT^{u+v}\)。
收缩映射可取为 \(\mathfrak X_0(M)\to \mathfrak X_U\) 的复合并限制到 \(U\)，于是 \(i\) 在 \(H^{1}\) 上诱导 \(\ZZ^{u}\hookrightarrow \ZZ^{u+v}\) 的直和因子嵌入。
由配对长正合列在低阶处的正合性，得到自然短正合列
\[
0\to H^{1}(\mathfrak X_0(M);\ZZ)\to H^{1}(\Sigma_0(M);\ZZ)\to H^{2}(\mathfrak X_0(M),\Sigma_0(M);\ZZ)\to 0,
\]
取秩可知 \(\rank H^{2}(\mathfrak X_0,\Sigma_0)=v\)。
并且在同构 \(H^{1}(\Sigma_0(M);\ZZ)\cong \Hom(G,\ZZ)\) 与
\(
H^{1}(\mathfrak X_0(M);\ZZ)\cong \Hom(U,\ZZ)
\)
之下，商群识别为 \(\Hom(G/U,\ZZ)\)，从而得到所述自然同构。
最后将两条秩等式代入定义 \ref{def:cdim-stokes-weighted-dimension}，得到 \(\mathrm{wdim}(M)=u+v/2\)；
而 \(\cdim^{\pm}(M)=\cdim(U)+\frac12\cdim(G/U)=u+\frac12 v\) 由定义 \ref{def:cdim-signed} 直接给出。证毕。
\end{proof}

\begin{corollary}[普适 Stokes 障碍格与连接同态的投影解释]\label{cor:cdim-stokes-universal-obstruction-lattice}
\leanverified{paper\_cdim\_stokes\_universal\_obstruction\_lattice}
在定理 \ref{thm:cdim-stokes-character-contraction-general-monoid} 的设定下，存在自然同构
\[
H^{2}\bigl(\mathfrak X_0(M),\Sigma_0(M);\ZZ\bigr)
\ \cong\
H^{1}\bigl(\Sigma_0(M);\ZZ\bigr)\big/ i^{\ast}H^{1}\bigl(\mathfrak X_0(M);\ZZ\bigr)
\ \cong\
\Hom(G/U,\ZZ).
\]
并且配对长正合列中的连接同态
\(
\delta:H^{1}(\Sigma_0(M);\ZZ)\to H^{2}(\mathfrak X_0(M),\Sigma_0(M);\ZZ)
\)
在上述同构下等同于对 \(\Hom(G,\ZZ)\) 的自然投影
\(
\Hom(G,\ZZ)\twoheadrightarrow \Hom(G/U,\ZZ)
\)。
因此 \(\Hom(G/U,\ZZ)\) 可视为由边界角色空间 \(\Sigma_0(M)\) 引出的普适 Stokes 障碍格：恰由那些在 \(\mathfrak X_0(M)\) 中可经模长塌缩到 \(0\) 而消去的相位方向生成。
\end{corollary}
\begin{proof}
由定理 \ref{thm:cdim-stokes-character-contraction-general-monoid} 证明中的短正合列立得第一同构。
而 \(H^{1}(\TT^{n};\ZZ)\cong \Hom(\ZZ^{n},\ZZ)\) 与 \(\Sigma_0(M)\cong \Hom(G,\TT)_0\) 给出 \(H^{1}(\Sigma_0(M);\ZZ)\cong \Hom(G,\ZZ)\)，同理 \(H^{1}(\mathfrak X_0(M);\ZZ)\cong \Hom(U,\ZZ)\)。
在该识别下，\(i^{\ast}\) 对应于 \(\Hom(G,\ZZ)\to \Hom(U,\ZZ)\) 的限制映射，故商群识别为 \(\Hom(G/U,\ZZ)\)，并且连接同态在正合性约束下等同于自然投影。证毕。
\end{proof}

\begin{theorem}[正交幺半群的 Stokes--上同调刻画]\label{thm:cdim-signed-stokes-cohomological-characterization}
设
\(
M=\ZZ^{u}\oplus \NN^{v}\oplus F
\)
（\(F\) 有限交换群）。
则
\[
\mathfrak X_0(M)\cong \TT^{u}\times \overline{\mathbb D}^{\,v},
\qquad
\Sigma_0(M)\cong \TT^{u+v},
\]
并且
\[
\boxed{\;\mathrm{wdim}(M)=\cdim^{\pm}(M)=u+\frac{v}{2}\;}
\qquad
(\text{与命题 \ref{prop:cdim-signed-orthant-closed} 一致}).
\]
\end{theorem}
\begin{proof}
由于幺半群同态必须将可逆元送入可逆元，有
\(
\Hom_{\mathrm{CMon}}(\ZZ,\overline{\mathbb D})=\TT
\)
而
\(
\Hom_{\mathrm{CMon}}(\NN,\overline{\mathbb D})=\overline{\mathbb D}
\)。
因此
\[
\mathfrak X(M)\cong \TT^{u}\times \overline{\mathbb D}^{\,v}\times \widehat F,
\qquad
\Sigma(M)\cong \TT^{u+v}\times \widehat F,
\]
其中 \(\widehat F\) 为有限离散空间；取包含平凡角色的连通分支即得第一行同胚。
\par
由于 \(\overline{\mathbb D}^{\,v}\) 可缩，
\(
H^{1}(\mathfrak X_0(M);\ZZ)\cong H^{1}(\TT^{u};\ZZ)\cong \ZZ^{u}
\)，故 \(\rank H^{1}(\mathfrak X_0(M))=u\)。
\par
再由配对 Künneth 公式与相对同调的悬挂同构
\(
H^{k}(\overline{\mathbb D}^{\,v},\TT^{v};\ZZ)\cong \widetilde H^{k-1}(\TT^{v};\ZZ)\ (k\ge 1)
\)，
可得
\(
H^{2}(\mathfrak X_0(M),\Sigma_0(M);\ZZ)\cong H^{2}(\overline{\mathbb D}^{\,v},\TT^{v};\ZZ)\cong H^{1}(\TT^{v};\ZZ)\cong \ZZ^{v}
\)，从而 \(\rank H^{2}(\mathfrak X_0(M),\Sigma_0(M))=v\)。
代入定义 \ref{def:cdim-stokes-weighted-dimension} 得 \(\mathrm{wdim}(M)=u+v/2\)，并与命题 \ref{prop:cdim-signed-orthant-closed} 的 \(\cdim^{\pm}(M)\) 一致。证毕。
\end{proof}
\leanverified{paper\_cdim\_signed\_stokes\_cohomological\_characterization}

\begin{proposition}[Stokes 连接同态与边界/体积的短正合列]\label{prop:cdim-stokes-exact-sequence-dictionary}
在定理 \ref{thm:cdim-signed-stokes-cohomological-characterization} 的设定下，
配对 \((\mathfrak X_0(M),\Sigma_0(M))\) 的长正合列在低维部分诱导短正合列
\[
0\longrightarrow H^{1}\bigl(\mathfrak X_0(M);\ZZ\bigr)
\longrightarrow H^{1}\bigl(\Sigma_0(M);\ZZ\bigr)
\xrightarrow{\ \delta\ }
H^{2}\bigl(\mathfrak X_0(M),\Sigma_0(M);\ZZ\bigr)
\longrightarrow 0.
\]
并且在自然同构
\(
H^{1}(\Sigma_0(M))\cong \ZZ^{u+v},\;
H^{1}(\mathfrak X_0(M))\cong \ZZ^{u},\;
H^{2}(\mathfrak X_0(M),\Sigma_0(M))\cong \ZZ^{v}
\)
之下，\(\delta\) 与坐标投影 \(\ZZ^{u+v}\twoheadrightarrow \ZZ^{v}\) 相合，核由前 \(u\) 个坐标张成。
\leanverified{paper\_cdim\_stokes\_exact\_sequence\_dictionary}
\end{proposition}
\begin{proof}
由
\(
\mathfrak X_0(M)\cong \TT^{u}\times \overline{\mathbb D}^{\,v}
\)
与
\(
\Sigma_0(M)=\TT^{u}\times \TT^{v}\hookrightarrow \TT^{u}\times \overline{\mathbb D}^{\,v}
\)
可知包含映射在 \(H^{1}\) 上诱导坐标嵌入 \(\ZZ^{u}\hookrightarrow \ZZ^{u+v}\)，其像为直和因子。
另一方面，由定理 \ref{thm:cdim-signed-stokes-cohomological-characterization} 的秩计算与正合性，连接同态 \(\delta\) 必满射且核恰为该嵌入像，从而得到短正合列。
\par
对最后陈述，取 \(\TT^{u+v}\) 的标准基 \(\{e_1,\dots,e_{u+v}\}\)（对应 \(H^{1}(\TT^{u+v};\ZZ)\cong \ZZ^{u+v}\) 的坐标基）。
前 \(u\) 个生成元来自 \(\TT^{u}\) 因子并在 \(\TT^{u}\times \overline{\mathbb D}^{\,v}\) 中延拓，因此落入 \(\ker\delta\)；
后 \(v\) 个生成元在每个圆盘对 \((\overline{\mathbb D},\TT)\) 上经连接同态送到相对基本类生成元，从而 \(\delta\) 与坐标投影一致。证毕。
\end{proof}

% Betti-average identity for wdim.

\begin{corollary}[边界--本体一阶 Betti 的平均律]\label{cor:cdim-wdim-betti-average-law}
令 \(M\) 为有限生成、可消去交换幺半群，并沿用定义 \ref{def:cdim-stokes-character-pair} 的 Stokes 角色空间配对
\(
(\mathfrak X_0(M),\Sigma_0(M))
\)
与定义 \ref{def:cdim-stokes-weighted-dimension} 的 \(\mathrm{wdim}(M)\)。
则成立恒等式
\[
\boxed{\;
\mathrm{wdim}(M)
=\frac12\Bigl(\rank_{\ZZ}H^{1}\bigl(\mathfrak X_0(M);\ZZ\bigr)+\rank_{\ZZ}H^{1}\bigl(\Sigma_0(M);\ZZ\bigr)\Bigr).\;}
\]
\leanverified{paper\_cdim\_wdim\_betti\_average\_law\_seeds}
\leanverified{paper\_cdim\_wdim\_betti\_average\_law}
\end{corollary}
\begin{proof}
由定理 \ref{thm:cdim-stokes-character-contraction-general-monoid} 的证明中低阶短正合列
\[
0\to H^{1}\bigl(\mathfrak X_0(M);\ZZ\bigr)\to H^{1}\bigl(\Sigma_0(M);\ZZ\bigr)\to H^{2}\bigl(\mathfrak X_0(M),\Sigma_0(M);\ZZ\bigr)\to 0,
\]
取秩得到
\(
\rank H^{2}(\mathfrak X_0,\Sigma_0)=\rank H^{1}(\Sigma_0)-\rank H^{1}(\mathfrak X_0)
\)。
代入定义 \ref{def:cdim-stokes-weighted-dimension} 即得所示平均律。证毕。
\end{proof}
\begin{proposition}[Betti 平均量的谱增长比较护栏]\label{prop:distill-connes-wdim-spectral-growth-guard}
令 \(M\) 满足推论 \ref{cor:cdim-wdim-betti-average-law} 的假设，并记
\[
w:=\mathrm{wdim}(M).
\]
设 \(D\) 为 Hilbert 空间上的正自伴紧预解算子，记其特征值计数函数为
\[
N_D(\Lambda):=\operatorname{Tr}\ind_{[0,\Lambda]}(D).
\]
若 \(w>0\)，且存在常数 \(C_-,C_+>0\) 与 \(\Lambda_0\ge1\)，使所有 \(\Lambda\ge\Lambda_0\) 满足
\[
C_-\Lambda^w\le N_D(\Lambda)\le C_+\Lambda^w,
\]
则
\[
\inf\{s>0:\operatorname{Tr}((1+D)^{-s})<\infty\}=w.
\]
另一方面，对任意 \(\alpha\ge0\)，存在正自伴紧预解算子 \(D_{\alpha}\)，使
\[
\inf\{s>0:\operatorname{Tr}((1+D_{\alpha})^{-s})<\infty\}=\alpha.
\]
因此，推论 \ref{cor:cdim-wdim-betti-average-law} 给出的 \(\mathrm{wdim}(M)\) 是一阶 Betti 平均读出；它等于某个算子谱维，必须由上述特征值计数或等价热迹比较作为额外账本保证。
\end{proposition}
\begin{proof}
先证比较命题。若 \(s>w\)，取 dyadic 壳层并用上界估计：
\[
\operatorname{Tr}((1+D)^{-s})\le C+\sum_{r\ge r_0}N_D(2^{r+1})2^{-rs}\le C+C_+2^w\sum_{r\ge r_0}2^{(w-s)r}<\infty.
\]
若 \(0<s<w\)，则对任意充分大的 \(\Lambda\)，部分和满足
\[
\operatorname{Tr}\bigl((1+D)^{-s}\ind_{[0,\Lambda]}(D)\bigr)
\ge (1+\Lambda)^{-s}N_D(\Lambda)
\ge c\Lambda^{w-s},
\]
其中 \(c>0\)。令 \(\Lambda\to\infty\) 得全迹发散。故临界指数为 \(w\)。

对任意 \(\alpha\ge0\)，在
\[
H_{\alpha}:=\bigoplus_{r=1}^{\infty}\CC^{m_r},\qquad
m_r:=\max\{1,\lceil 2^{\alpha r}\rceil\}
\]
上定义 \(D_{\alpha}|_{\CC^{m_r}}=2^r\mathrm{Id}\)。则
\[
\operatorname{Tr}((1+D_{\alpha})^{-s})\asymp
\sum_{r\ge1}m_r2^{-rs}.
\]
与前一证明同样的几何级数比较给出临界指数 \(\alpha\)；当 \(\alpha=0\) 时该级数对所有 \(s>0\) 收敛。证毕。
\end{proof}
\begin{definition}[边界相对 Stokes--Euler 路径账本]\label{def:distill-ecalle-relative-stokes-euler-ledger}
在推论 \ref{cor:cdim-wdim-betti-average-law} 的假设下，记
\[
A:=H^{1}\bigl(\mathfrak X_0(M);\ZZ\bigr),\qquad
B:=H^{1}\bigl(\Sigma_0(M);\ZZ\bigr),\qquad
C:=H^{2}\bigl(\mathfrak X_0(M),\Sigma_0(M);\ZZ\bigr).
\]
并令
\[
0\to A\to B\xrightarrow{\pi} C\to0
\]
为该推论证明中使用的低阶短正合列。一个边界相对 Stokes--Euler 路径账本是四元组
\[
(\mathfrak S,c_\sigma,u_\sigma,\nu),
\]
其中 \(\mathfrak S\) 是有限地址集，\(c_\sigma\in C\) 是地址 \(\sigma\) 的相对 Stokes 跃迁类，\(u_\sigma\) 属于某个自由阿贝尔群 \(U\)，而
\(\nu:U\to\ZZ_{\ge0}\) 是满足
\(\nu(u+u')\le\nu(u)+\nu(u')\) 的 Euler 尾项预算函数。对地址字
\(\gamma=(\sigma_1,\ldots,\sigma_k)\)，定义
\[
c(\gamma):=c_{\sigma_1}+\cdots+c_{\sigma_k},\qquad
u(\gamma):=\nu(u_{\sigma_1}+\cdots+u_{\sigma_k}).
\]
若 \(\Gamma\) 是有限路径族，则称
\[
J_\Gamma:=\left\langle c(\gamma)-c(\gamma'):\gamma,\gamma'\in\Gamma
\text{ 且端点相同}\right\rangle_{\ZZ}\subseteq C
\]
为该路径族的相对 Stokes holonomy 模块。
\end{definition}

\begin{proposition}[相对 Stokes 路径和的 Betti 平均审计]\label{prop:distill-ecalle-relative-stokes-euler-betti-audit}
沿用定义 \ref{def:distill-ecalle-relative-stokes-euler-ledger}。对任意子群
\(E\le C\)，令
\[
B_E:=\pi^{-1}(E)\subseteq B,\qquad
\mathrm{wdim}_E(M):=\frac12\bigl(\rank_{\ZZ}A+\rank_{\ZZ}B_E\bigr).
\]
则
\[
\mathrm{wdim}_E(M)=\rank_{\ZZ}A+\frac12\rank_{\ZZ}E.
\]
特别地，对有限路径族 \(\Gamma\)，相对 Stokes holonomy 模块给出的边界平均偏差为
\[
\mathrm{wdim}_{J_\Gamma}(M)-\mathrm{wdim}_{0}(M)
=\frac12\rank_{\ZZ}J_\Gamma.
\]
若对该账本还满足预算支配条件
\[
\rank_{\ZZ}\langle c_\sigma:\sigma\in\mathfrak S\rangle_{\ZZ}\le
\sum_{\sigma\in\mathfrak S}\nu(u_\sigma),
\]
则对任意有限路径族 \(\Gamma\) 有
\[
\rank_{\ZZ}J_\Gamma\le
\sum_{\sigma\in\mathfrak S}\nu(u_\sigma).
\]
因此，边界 Stokes 路径和只有在 \(J_\Gamma=0\) 时才是相对同调意义下路径无关；若
\(J_\Gamma\neq0\)，其秩给出必须被边界残差账本记录的 holonomy 宽度。
\end{proposition}
\begin{proof}
由于
\(0\to A\to B\xrightarrow{\pi}C\to0\) 正合，限制到 \(B_E=\pi^{-1}(E)\) 得到短正合列
\[
0\to A\to B_E\to E\to0.
\]
取有限生成阿贝尔群的秩，得
\(\rank B_E=\rank A+\rank E\)，从而
\[
\mathrm{wdim}_E(M)=\frac12(\rank A+\rank A+\rank E)
=\rank A+\frac12\rank E.
\]
令 \(E=J_\Gamma\) 即得第二个等式。

地址集 \(\mathfrak S\) 把边界跃迁分解为局部相对类 \(c_\sigma\)。任一路径差
\(c(\gamma)-c(\gamma')\) 都属于由 \(c_\sigma\) 生成的子群，因此
\[
J_\Gamma\le \langle c_\sigma:\sigma\in\mathfrak S\rangle_{\ZZ}.
\]
预算支配条件随即给出所需秩界。最后，\(J_\Gamma=0\) 等价于任意同端点路径的相对 Stokes 和相等；若该模块非零，则存在同端点路径差在相对群 \(C\) 中非平凡，不能由边界一阶读出自动消去，只能作为有限地址支撑上的 holonomy 宽度记录。
\end{proof}

\endinput

% Amalgamation along units for orthogonal monoids and the induced wdim formula.

\begin{theorem}[正交幺半群沿单位子幺半群并合的 Stokes 维数公式]\label{thm:cdim-orthogonal-amalgamated-wdim}
设 \(M_1,M_2\) 为正交幺半群，并存在同一可逆子幺半群 \(U\) 及单态
\(
\iota_i:U\hookrightarrow U(M_i)\subseteq M_i
\)
（\(i=1,2\)）。
令 \(M\) 为交换幺半群范畴中由 \(\iota_1,\iota_2\) 诱导的并合（推送出）：
\[
M:=M_1\amalg_{U}M_2.
\]
记 \(U_{\mathrm{free}}\) 为 \(U\) 的自由阿贝尔部分，则 \(M\) 仍为正交幺半群，并且满足
\[
\boxed{\;
\mathrm{wdim}(M)=\mathrm{wdim}(M_1)+\mathrm{wdim}(M_2)-\rank_{\ZZ}U_{\mathrm{free}}.\;}
\]
\end{theorem}
\begin{proof}
由正交假设，存在分解
\[
M_i\cong \ZZ^{u_i}\oplus \NN^{v_i}\oplus F_i
\qquad(i=1,2),
\]
其中 \(F_i\) 为有限交换群。
因此其单位群满足
\(
U(M_i)\cong \ZZ^{u_i}\oplus F_i
\)，并且群完成满足
\(
G(M_i)\cong \ZZ^{u_i+v_i}\oplus F_i
\)。
\par
由于 \(\iota_i(U)\subseteq U(M_i)\)，并合仅在单位方向施加识别关系；而 \(\NN^{v_1}\) 与 \(\NN^{v_2}\) 的正锥部分在并合中保持直和（否则将迫使非单位元素在并合中可逆，与正交结构矛盾）。
因此 \(G(M)\) 的无挠秩满足
\[
\rank_{\ZZ}G(M)_{\mathrm{free}}=(u_1+v_1)+(u_2+v_2)-\rank_{\ZZ}U_{\mathrm{free}},
\]
而单位群的无挠秩满足
\[
\rank_{\ZZ}U(M)_{\mathrm{free}}=u_1+u_2-\rank_{\ZZ}U_{\mathrm{free}}.
\]
从而
\[
\rank_{\ZZ}\bigl(G(M)/U(M)\bigr)=v_1+v_2.
\]
由定理 \ref{thm:cdim-stokes-character-contraction-general-monoid} 的恒等式
\(
\mathrm{wdim}(N)=\rank_{\ZZ}U(N)+\frac12\rank_{\ZZ}(G(N)/U(N))
\)
对 \(N=M,M_1,M_2\) 分别应用并代入上述秩计算，即得
\[
\mathrm{wdim}(M)
=\Bigl(u_1+u_2-\rank U_{\mathrm{free}}\Bigr)+\frac12(v_1+v_2)
=\mathrm{wdim}(M_1)+\mathrm{wdim}(M_2)-\rank U_{\mathrm{free}}.
\]
证毕。
\end{proof}

\endinput


% Homology version of the Stokes exact sequence and splitting.

\begin{proposition}[Stokes 连接同态的同调短正合列与单位/非单位分裂]\label{prop:cdim-stokes-homology-exact-splitting}
\leanverified{paper\_cdim\_stokes\_homology\_exact\_splitting}
沿用定义 \ref{def:cdim-stokes-character-pair} 的记号，令包含映射 \(i:\Sigma_0(M)\hookrightarrow \mathfrak X_0(M)\)。
则配对 \((\mathfrak X_0(M),\Sigma_0(M))\) 的长正合列在低维同调上诱导短正合列
\[
0\longrightarrow H_{2}\bigl(\mathfrak X_0(M),\Sigma_0(M);\ZZ\bigr)
\xrightarrow{\ \partial\ }
H_{1}\bigl(\Sigma_0(M);\ZZ\bigr)
\xrightarrow{\ i_\ast\ }
H_{1}\bigl(\mathfrak X_0(M);\ZZ\bigr)
\longrightarrow 0,
\]
其中 \(\partial\) 为相对边界映射。
并且在定理 \ref{thm:cdim-stokes-character-contraction-general-monoid} 的识别下，存在基选择使得
\[
H_{1}\bigl(\Sigma_0(M);\ZZ\bigr)\cong \ZZ^{u}\oplus \ZZ^{v},\qquad
H_{1}\bigl(\mathfrak X_0(M);\ZZ\bigr)\cong \ZZ^{u},
\]
且 \(i_\ast\) 在该分解下等同于投影到前 \(u\) 个坐标，而 \(\partial\) 将
\(
H_{2}(\mathfrak X_0(M),\Sigma_0(M);\ZZ)\cong \ZZ^{v}
\)
同构到第二个直和因子。
\end{proposition}
\begin{proof}
由定理 \ref{thm:cdim-stokes-character-contraction-general-monoid}，
\(
\mathfrak X_0(M)\simeq \TT^{u}
\)、
\(
\Sigma_0(M)\cong \TT^{u+v}
\)，
并且包含映射在同伦意义下等同于 \(\TT^{u+v}\to \TT^{u}\) 的坐标投影。
因此 \(i_\ast:H_1(\Sigma_0;\ZZ)\to H_1(\mathfrak X_0;\ZZ)\) 为满射，核为自由秩 \(v\)。
\par
另一方面，\(H_2(\mathfrak X_0(M);\ZZ)\cong \wedge^2\ZZ^{u}\)、\(H_2(\Sigma_0(M);\ZZ)\cong \wedge^2\ZZ^{u+v}\)，
而由 \(H_1\) 上的满射诱导 \(\wedge^2\ZZ^{u+v}\twoheadrightarrow \wedge^2\ZZ^{u}\)，
故 \(H_2(\Sigma_0)\to H_2(\mathfrak X_0)\) 亦满射。
将此代入同调长正合列片段
\[
H_2(\Sigma_0)\to H_2(\mathfrak X_0)\to H_2(\mathfrak X_0,\Sigma_0)\xrightarrow{\partial}H_1(\Sigma_0)\xrightarrow{i_\ast}H_1(\mathfrak X_0)\to 0
\]
得到
\(
0\to H_2(\mathfrak X_0,\Sigma_0)\xrightarrow{\partial}\ker i_\ast\to 0
\)，
即 \(H_2(\mathfrak X_0,\Sigma_0)\cong \ker i_\ast\) 且由 \(\partial\) 实现嵌入，从而得到所示短正合列。
\par
最后，对自由阿贝尔群可取适配基使满射 \(i_\ast\) 呈块矩阵 \((I_u\ \ 0)\)，此时 \(\ker i_\ast\cong \ZZ^{v}\) 由后 \(v\) 个坐标生成，而 \(\partial\) 给出同构到该直和因子。证毕。
\end{proof}

\endinput

% Conjectural unification of Walsh--Stokes readout with Stokes character spaces.

\begin{conjecture}[Stokes 角色空间与 Walsh--Stokes 字典的统一同调解释]\label{conj:cdim-walsh-stokes-unified-homology}
对每个整数 \(n\ge 1\) 与任意满射粗粒化 \(f:\Omega_n=\{0,1\}^n\twoheadrightarrow X\)（见 \S\ref{subsubsec:spg-boundary-affine-residual-freedom}），存在一个由 \((\Omega_n,f)\) 自然诱导的有限生成、可消去交换幺半群 \(M_{n,f}\)，满足：
\begin{enumerate}
  \item \textbf{（边界回路秩的角色化）}
  边界邻接多图 \(B_f\)（定义于 \S\ref{subsubsec:spg-boundary-affine-residual-freedom}）的第一同调与边界角色空间的第一同调同构：
  \[
  H_1(B_f;\ZZ)\ \cong\ H_1\bigl(\Sigma_0(M_{n,f});\ZZ\bigr).
  \]
  \item \textbf{（Walsh 读出之连接因子化）}
  对任意指标集 \(I\subseteq\{1,\dots,n\}\)，定理 \ref{thm:spg-walsh-discrete-stokes-holography} 的 \(I\)-阶差分读出
  \(
  f\mapsto b_I(f)
  \)
  可经由配对 \((\mathfrak X_0(M_{n,f}),\Sigma_0(M_{n,f}))\) 的相对同调连接同态（命题 \ref{prop:cdim-stokes-homology-exact-splitting}）因子化；
  换言之，Walsh--Stokes 读出与相对链的边界/体积转移在同调意义下同构。
  \item \textbf{（半整数计数与外置自由度）}
  由定义 \ref{def:cdim-stokes-weighted-dimension} 所给出的半整数 \(\mathrm{wdim}(M_{n,f})\) 与边界回路秩 \(r(f)=\rank_{\ZZ}H_1(B_f;\ZZ)\)（定理 \ref{thm:spg-boundary-affine-fiber-cycle-rank}）之间满足可审计的计数一致性：其“非单位方向”贡献精确反映 \(r(f)\) 所控制的边界剩余自由度规模。
\end{enumerate}
\end{conjecture}

\endinput

% De Rham lift of discrete Stokes on the hypercube via Sigma_0(N^n) = T^n.

\begin{proposition}[\texorpdfstring{$\Sigma_0(\NN^n)=\TT^n$}{Sigma0(N^n)=T^n} 上的 de Rham 边界读出：离散 Stokes 的角色空间提升]\label{prop:cdim-hypercube-derham-stokes-lift}
令 \(n\ge 1\)，\(\Omega_n:=\{0,1\}^n\)，并将
\(
\Sigma_0(\NN^n)\cong \Hom(\ZZ^n,\TT)_0
\)
自然识别为 \(n\)-维环面 \(\TT^n\)，其角坐标记为 \((\theta_1,\dots,\theta_n)\)。
定义 \(2\)-挠点嵌入
\[
\iota:\Omega_n\hookrightarrow \TT^n,
\qquad
\iota(\omega):=(e^{\pi i\omega_1},\dots,e^{\pi i\omega_n})\in\{\pm 1\}^n\subset\TT^n.
\]
对超立方体有向边 \(e:\omega\to\omega^{(i)}\)（翻转第 \(i\) 位）定义分段光滑路径 \(\gamma_e:[0,1]\to\TT^n\)：
除第 \(i\) 个坐标外其余坐标保持常数，仅令 \(\theta_i\) 以常速沿长度为 \(\pi\) 的最短弧从 \(\pi\omega_i\) 变到 \(\pi(1-\omega_i)\)。
令 \(S\subseteq\Omega_n\)，并记 \(\partial^{\to}S\) 为从 \(S\) 指向补集的外法向有向边集合。
则对每个坐标 \(i\) 成立精确恒等式
\[
\boxed{\;
b_i(S)=\frac{1}{\pi}\sum_{e\in\partial^{\to}S}\int_{\gamma_e}\dd\theta_i\;}
\]
其中 \(b_i(S)=\sum_{\omega\in S}(-1)^{\omega_i}\) 为第 \(i\) 坐标体积偏差（命题 \ref{prop:fold-hypercube-godel-stokes-flux-bias}）。
\leanverified{pathIntegralCoeff}
\leanverified{pathIntegralCoeff\_false}
\leanverified{pathIntegralCoeff\_true}
\leanverified{pathIntegralCoeff\_abs}
\leanverified{pathIntegralCoeff\_sq}
\leanverified{pathIntegral}
\leanverified{pathIntegral\_false}
\leanverified{pathIntegral\_true}
\leanverified{pathIntegral\_div\_pi}
\leanverified{paper\_cdim\_hypercube\_derham\_stokes\_lift}
\end{proposition}
\begin{proof}
按构造，若 \(e:\omega\to\omega^{(i)}\)，则对 \(j\neq i\) 有 \(\int_{\gamma_e}\dd\theta_j=0\)，并且
\(
\int_{\gamma_e}\dd\theta_i=\pi
\)
当 \(\omega_i=0\)，而
\(
\int_{\gamma_e}\dd\theta_i=-\pi
\)
当 \(\omega_i=1\)。
因此
\[
\sum_{e\in\partial^{\to}S}\int_{\gamma_e}\dd\theta_i
=\pi\bigl(N_i^+(S)-N_i^-(S)\bigr),
\]
其中 \(N_i^\pm(S)\) 为命题 \ref{prop:fold-hypercube-godel-stokes-flux-bias} 的有向外法边计数。
再由同一命题的离散 Stokes 恒等式 \(N_i^+(S)-N_i^-(S)=b_i(S)\) 得证。证毕。
\end{proof}

\endinput



\begin{proposition}[相对上同调环的分解与高阶 Stokes 结构]\label{prop:cdim-stokes-relative-cohomology-ring}
\leanverified{paper\_cdim\_stokes\_relative\_cohomology\_ring}
在定理 \ref{thm:cdim-signed-stokes-cohomological-characterization} 的设定下，有自然同构（作为分次交换环）
\[
H^{*}\bigl(\mathfrak X_0(M),\Sigma_0(M);\ZZ\bigr)
\ \cong\
H^{*}(\TT^{u};\ZZ)\otimes \widetilde H^{*-1}(\TT^{v};\ZZ),
\]
从而对每个 \(k\ge 1\) 有分解
\[
H^{k}\bigl(\mathfrak X_0(M),\Sigma_0(M);\ZZ\bigr)
\ \cong\
\bigoplus_{p+q=k}\Bigl(\wedge^{p}\ZZ^{u}\Bigr)\otimes\Bigl(\wedge^{q-1}\ZZ^{v}\Bigr),
\]
并且右端楔代数乘法与 cup-product 相容。
\end{proposition}
\begin{proof}
由配对 Künneth 公式，
\(
H^{*}(\TT^{u}\times \overline{\mathbb D}^{\,v},\TT^{u}\times \TT^{v})
\cong
H^{*}(\TT^{u})\otimes H^{*}(\overline{\mathbb D}^{\,v},\TT^{v})
\)。
另一方面，由 \(\overline{\mathbb D}^{\,v}\) 可缩的长正合列可得对 \(k\ge 1\)
\(
H^{k}(\overline{\mathbb D}^{\,v},\TT^{v})\cong \widetilde H^{k-1}(\TT^{v})
\)。
再用 \(H^{*}(\TT^{u};\ZZ)\cong \wedge^{*}\ZZ^{u}\)、\(\widetilde H^{*}(\TT^{v};\ZZ)\cong \wedge^{*}\ZZ^{v}\) 即得。证毕。
\end{proof}

\begin{theorem}[Stokes 转迁的显式相对上同调模型与可计算基]\label{thm:cdim-stokes-explicit-relative-cohomology-basis}
\leanverified{paper\_cdim\_stokes\_explicit\_relative\_cohomology\_basis}
设 \(v\ge 1\)，并将闭单位圆盘 \(\overline{\mathbb D}\) 视为带角流形。
在 \(\overline{\mathbb D}^{\,v}\) 上取极坐标 \(z_j=r_j e^{\mathrm{i}\phi_j}\)（\(1\le j\le v\)）。
对任意非空指标集 \(I\subseteq[v]\) 定义微分形式
\[
\alpha_{I}:=\bigwedge_{j\in I}\left(\frac{r_j^2}{2\pi}\,d\phi_j\right)\in\Omega^{|I|}(\overline{\mathbb D}^{\,v}),
\qquad
\Omega_{I}:=d\alpha_I\in\Omega^{|I|+1}(\overline{\mathbb D}^{\,v}).
\]
则：
\begin{enumerate}
  \item \(\Omega_I|_{\TT^{v}}=0\)，从而 \(\Omega_I\) 定义一个相对闭类
  \(
  [\Omega_I]\in H^{|I|+1}(\overline{\mathbb D}^{\,v},\TT^{v};\ZZ)
  \)；
  并且当 \(|I|=k\) 时，\(\{[\Omega_I]:|I|=k\}\) 构成 \(H^{k+1}(\overline{\mathbb D}^{\,v},\TT^{v};\ZZ)\) 的一组 \(\ZZ\)-基。
  \item 设 \(\delta:H^{k}(\TT^{v};\ZZ)\to H^{k+1}(\overline{\mathbb D}^{\,v},\TT^{v};\ZZ)\) 为配对长正合列的连接同态。
  则在 \(\TT^{v}\) 的标准外代数基
  \(
  \bigwedge_{j\in I}\frac{d\phi_j}{2\pi}\in H^{k}(\TT^{v};\ZZ)
  \)
  上有显式公式
  \[
  \boxed{\;
  \delta\!\left(\bigwedge_{j\in I}\frac{d\phi_j}{2\pi}\right)=[\Omega_I]\; } .
  \]
  特别地，\(\delta:H^{1}(\TT^{v};\ZZ)\xrightarrow{\sim}H^{2}(\overline{\mathbb D}^{\,v},\TT^{v};\ZZ)\) 为同构，
  且与命题 \ref{prop:cdim-stokes-exact-sequence-dictionary} 的坐标化描述一致。
  \item 对任意 \(u\ge 0\)，在识别
  \(
  (\mathfrak X_0(\ZZ^{u}\oplus\NN^{v}),\Sigma_0(\ZZ^{u}\oplus\NN^{v}))
  \cong
  (\TT^{u}\times \overline{\mathbb D}^{\,v},\TT^{u}\times\TT^{v})
  \)
  下，右乘 cup-product 与 Künneth 分解给出
  \[
  H^{m}\bigl(\mathfrak X_0,\Sigma_0;\ZZ\bigr)
  \cong
  \bigoplus_{p+q=m}H^{p}(\TT^{u};\ZZ)\otimes H^{q}(\overline{\mathbb D}^{\,v},\TT^{v};\ZZ),
  \]
  从而由 \(\TT^u\) 的标准外代数基 \([\eta]\) 与 \([\Omega_I]\) 的 cup-product
  \(
  [\eta]\smile[\Omega_I]
  \)
  得到一组显式可计算基；该基与命题 \ref{prop:cdim-stokes-relative-cohomology-ring} 的分次环分解相容。
\end{enumerate}
\end{theorem}
\begin{proof}
对任意 \(I\)，由于在边界 \(\TT^{v}\subset \overline{\mathbb D}^{\,v}\) 上各 \(r_j\equiv 1\) 为常数，故 \(dr_j|_{\TT^{v}}=0\)，从而 \(\Omega_I=d\alpha_I\) 限制到边界为零。
因此 \(\Omega_I\) 代表相对闭类。
\par
另一方面，\(\alpha_I|_{\TT^{v}}=\bigwedge_{j\in I}\frac{d\phi_j}{2\pi}\)。
由相对 de Rham 模型中连接同态的链级刻画，若 \(\beta\in\Omega^{k}(\overline{\mathbb D}^{\,v})\) 满足 \(\beta|_{\TT^v}\) 代表某个边界上同调类，则 \(\delta([\beta|_{\TT^v}])=[d\beta]\)。
取 \(\beta=\alpha_I\) 即得 \(\delta(\bigwedge_{j\in I}\frac{d\phi_j}{2\pi})=[\Omega_I]\)。
\par
由于 \(\overline{\mathbb D}^{\,v}\) 可缩，长正合列给出对 \(k\ge 1\) 的同构
\(
H^{k+1}(\overline{\mathbb D}^{\,v},\TT^{v};\ZZ)\cong H^{k}(\TT^{v};\ZZ)
\)，且由 \(\delta\) 实现；结合上式可知 \(\{[\Omega_I]:|I|=k\}\) 与 \(\TT^{v}\) 的标准外代数基一一对应，故构成 \(\ZZ\)-基。
\par
最后一条由配对 Künneth 公式与 \([\Omega_I]\) 的基性质直接推出，并与命题 \ref{prop:cdim-stokes-relative-cohomology-ring} 的环同构相容。证毕。
\end{proof}

\begin{theorem}[相对第一陈类与边界框架分类的 Stokes 同构]\label{thm:cdim-stokes-relative-c1-framed-line-bundles}
令 \(n\ge 1\)。设 \(\mathcal L^{\mathrm{fr}}(\overline{\mathbb D}^{\,n})\) 为带边界框架的复线丛同构类集合：
其元素为对 \((L,\tau)\)，其中 \(L\to \overline{\mathbb D}^{\,n}\) 为复线丛，\(\tau\) 为其在边界 \(\partial\overline{\mathbb D}^{\,n}=\TT^{n}\) 上的一个平凡化。
则存在自然同构
\[
\boxed{\;
c_1^{\mathrm{rel}}:\ \mathcal L^{\mathrm{fr}}(\overline{\mathbb D}^{\,n})\xrightarrow{\ \sim\ }H^{2}(\overline{\mathbb D}^{\,n},\TT^{n};\ZZ)\; } .
\]
并且当 \(L\) 为平凡丛 \(L=\overline{\mathbb D}^{\,n}\times\CC\) 且框架 \(\tau\) 由夹持函数 \(g_\tau:\TT^{n}\to U(1)\) 给出时，有
\[
\boxed{\;
c_1^{\mathrm{rel}}(L,\tau)=\delta([g_\tau])\in H^{2}(\overline{\mathbb D}^{\,n},\TT^{n};\ZZ)\; } ,
\]
其中 \(\delta:H^{1}(\TT^{n};\ZZ)\to H^{2}(\overline{\mathbb D}^{\,n},\TT^{n};\ZZ)\) 为连接同态（见定理 \ref{thm:cdim-stokes-explicit-relative-cohomology-basis} 与命题 \ref{prop:cdim-stokes-exact-sequence-dictionary}）。
此外，\(c_1^{\mathrm{rel}}\) 可由 Stokes 配对刻画为“边界绕数 \(=\) 体内曲率积分”：若取任意 unitary 连接 \(A\) 使其边界规范势与 \(\tau\) 相容，令 \(F=dA\)，则对任意相对 \(2\)-链 \(C\) 有
\[
\int_{C}F=\int_{\partial C}A.
\]
\leanverified{paper\_cdim\_stokes\_relative\_c1\_framed\_line\_bundles}
\end{theorem}
\begin{proof}
\(\overline{\mathbb D}^{\,n}\) 可缩，故任意复线丛在体内平凡；带边界框架的同构类因此等价于边界过渡函数的同伦类
\(
[\TT^{n},U(1)]\cong H^{1}(\TT^{n};\ZZ)
\)。
另一方面，由连接同态 \(\delta:H^{1}(\TT^{n};\ZZ)\xrightarrow{\sim}H^{2}(\overline{\mathbb D}^{\,n},\TT^{n};\ZZ)\)（定理 \ref{thm:cdim-stokes-explicit-relative-cohomology-basis}）得到自然同构，定义为 \(c_1^{\mathrm{rel}}\)。
\par
当 \(L\) 为平凡丛时，\(\tau\) 的同伦类由 \(g_\tau\) 给出，且上同构在该表示下即为 \(\delta([g_\tau])\)。
最后的 Stokes 配对刻画由 \(F=dA\) 与带边界的 Stokes 公式（命题 \ref{prop:cdim-stokes-half-face-count} 的局部模型）直接推出。证毕。
\end{proof}

\begin{proposition}[\texorpdfstring{$q$}{q}-幂映射的 Stokes 缩放律与异常放大的上同调实现]\label{prop:cdim-stokes-qpower-scaling-anom-amplification}
令 \(n\ge 1\)，并定义坐标 \texorpdfstring{$q$}{q}-幂映射
\[
P_q:\overline{\mathbb D}^{\,n}\to \overline{\mathbb D}^{\,n},
\qquad
(z_1,\dots,z_n)\mapsto (z_1^{q},\dots,z_n^{q})
\qquad(q\in\NN_{\ge 1}).
\]
则：
\begin{enumerate}
  \item 在相对上同调上，拉回
  \(
  P_q^\ast:H^{2}(\overline{\mathbb D}^{\,n},\TT^{n};\ZZ)\to H^{2}(\overline{\mathbb D}^{\,n},\TT^{n};\ZZ)
  \)
  等于标量乘法 \(q\)；
  更一般地，对任意非空 \(I\subseteq[n]\) 有
  \[
  \boxed{\;P_q^\ast[\Omega_I]=q^{|I|}\,[\Omega_I]\;}
  \quad\text{于}\quad
  H^{|I|+1}(\overline{\mathbb D}^{\,n},\TT^{n};\ZZ),
  \]
  其中 \([\Omega_I]\) 为定理 \ref{thm:cdim-stokes-explicit-relative-cohomology-basis} 的基类。
  \item 设 \(G\) 为有限阿贝尔群，并记异常签名向量空间
  \(
  \mathcal{A}_G:=\RR^{\widehat G\setminus\{\chi_0\}}
  \)
  （定义 \ref{def:pom-anom-accumulation}）。
  令 \(M_G:=\NN^{\widehat G\setminus\{\chi_0\}}\) 为自由正交交换幺半群，则
  \[
  (\mathfrak X_0(M_G),\Sigma_0(M_G))\cong (\overline{\mathbb D}^{\,n},\TT^{n}),\qquad n:=|\widehat G|-1,
  \]
  并存在一个按坐标基识别的自然线性同构
  \[
  \boxed{\;
  \mathcal{A}_G\ \cong\ H^{2}\bigl(\mathfrak X_0(M_G),\Sigma_0(M_G);\RR\bigr)\;}
  \]
  使得命题 \ref{prop:pom-anom-gap-amplification} 的放大律
  \(
  [\mathrm{Anom}(w^{\boxtimes q})]=q[\mathrm{Anom}(w)]
  \)
  与 \(P_q^\ast\) 的上同调缩放在该模型下对齐。
\end{enumerate}
\end{proposition}
\begin{proof}
(1) 在边界 \(\TT^{n}\) 上，\(P_q|_{\TT^{n}}\) 在每个圆周因子上为度数 \(q\) 的覆盖，故在 \(H^{1}(\TT^{n};\ZZ)\) 上为乘以 \(q\)，并在 \(\wedge^{|I|}\) 上为乘以 \(q^{|I|}\)。
由连接同态 \(\delta\) 的自然性与定理 \ref{thm:cdim-stokes-explicit-relative-cohomology-basis}(2) 的识别
\(
[\Omega_I]=\delta(\bigwedge_{j\in I}\frac{d\phi_j}{2\pi})
\)
立得 \(P_q^\ast[\Omega_I]=q^{|I|}[\Omega_I]\)；
取 \(|I|=1\) 即得相对 \(H^2\) 上的标量乘法 \(q\)。
\par
(2) 对 \(M_G=\NN^{n}\)，由定义 \ref{def:cdim-stokes-character-pair} 有
\(
\mathfrak X_0(M_G)\cong \overline{\mathbb D}^{\,n}
\)、
\(
\Sigma_0(M_G)\cong \TT^{n}
\)。
而 \(H^{2}(\overline{\mathbb D}^{\,n},\TT^{n};\RR)\cong \RR^{n}\) 的坐标基可取为 \(\{[\Omega_{\{j\}}]\}\)；
与 \(\mathcal{A}_G=\RR^{n}\) 的标准基对齐即可得到自然同构。
在该识别下，“放大律为乘以 \(q\)”与第 (1) 条在相对 \(H^2\) 上的缩放一致。证毕。
\end{proof}
\leanverified{paper\_cdim\_stokes\_qpower\_scaling\_anom\_amplification}

\endinput


\begin{proposition}[单侧角流形的 Stokes 面计数与半权重]\label{prop:cdim-stokes-half-face-count}
令
\(
\Omega_{a,b}:=\RR^{a}\times[0,\infty)^{b}
\)
取标准取向。
对任意紧支撑光滑 \((a{+}b{-}1)\)-形式 \(\omega\) 于 \(\Omega_{a,b}\)，广义 Stokes 定理给出
\[
\int_{\Omega_{a,b}} d\omega
\;=\;
\sum_{j=1}^{b}\pm\int_{F_j}\omega|_{F_j},
\qquad
F_j:=\RR^{a}\times\{x_j=0\}\times[0,\infty)^{b-1},
\]
并且由于紧支撑假设，无穷远端不产生边界项。
与之对比，对紧超立方体 \(C_{a,b}:=\RR^{a}\times[0,1]^{b}\)，边界由 \(2b\) 个互不相交面构成；
若对同类 \(\omega\) 施加 Dirichlet 条件使其在外侧面 \(\{x_j=1\}\) 上为零，则 Stokes 边界项仅保留 \(b\) 个内侧面 \(\{x_j=0\}\)。
因此单侧通道在 Stokes 边界计数中仅贡献一个面，而两侧可逆通道在紧化后贡献两个面；
该折半机制与命题 \ref{prop:cdim-signed-orthant-closed} 的 \(u+\tfrac v2\) 权重完全对齐。
\end{proposition}
\leanverified{paper\_cdim\_stokes\_half\_face\_count}
\begin{proof}
第一式为带角流形上的 Stokes 公式；紧支撑保证无穷远端边界积分为零。
第二段为同一公式在 \(C_{a,b}\) 上对各个面逐一展开，并在 Dirichlet 条件下消去外侧面贡献。
最后陈述由“单侧面数折半、两侧面数成对”的计数与命题 \ref{prop:cdim-signed-orthant-closed} 的权重一致性直接推出。证毕。
\end{proof}

\begin{proposition}[角流形纤维积分的 Stokes 缺陷项与单侧面贡献机制]\label{prop:cdim-stokes-fiber-integration-defect}
设 \(\pi:E\to B\) 为定向、适当的光滑子沉，允许角结构，纤维为紧致 \(d\)-维流形。
记水平边界 \(\partial_hE:=\pi^{-1}(\partial B)\) 与垂直边界 \(\partial_vE\)（由纤维边界/塌缩面诱导的边界分量），并令 \(\pi^{\partial}:\partial_vE\to B\) 为诱导映射。
对相对形式复形
\[
\Omega_{\mathrm{rel}}^{k}(E):=\{\omega\in\Omega^{k}(E):\ \omega|_{\partial_hE}=0\}
\]
定义纤维积分算子 \(\pi_!:\Omega_{\mathrm{rel}}^{k}(E)\to \Omega^{k-d}(B)\)。
则对任意 \(\omega\in\Omega_{\mathrm{rel}}^{k}(E)\) 有精确恒等式
\[
\boxed{\;
d(\pi_!\omega)=(-1)^d\,\pi_!(d\omega)\;+\;(\pi^{\partial})_!\bigl(\omega|_{\partial_vE}\bigr)\; } .
\]
特别地：
\begin{enumerate}
  \item 若进一步施加 Dirichlet 型消失条件使 \(\omega\) 在 \(\partial_vE\) 的某些面分量上亦为零，则边界项仅由其余面分量贡献；
  这正是命题 \ref{prop:cdim-stokes-half-face-count} 所强调的“单侧面贡献/半权重”机制在纤维推前口径下的链级实现。
  \item 将该恒等式应用于命题 \ref{prop:cdim-stokes-radial-fibration-transgression} 的径向塌缩映射
  \(
  \rho:\mathfrak X_0(M)\to \TT^u\times[0,1]^v
  \)
  时，\((\rho^{\partial})_!(\omega|_{\partial_v})\) 的同调类落入由定理 \ref{thm:cdim-stokes-explicit-relative-cohomology-basis} 的相对基类 \([\Omega_{\{j\}}]\) 张成的子空间；
  因而 Beck--Chevalley 交换式的缺陷在该模型中被一个显式可基表示的相对 \(H^2\)-残差完全控制。
\end{enumerate}
\end{proposition}
\begin{proof}
在每点的局部坐标中，\(E\) 可表示为 \(\Omega_{a,b}\times F\) 的适当开集（其中 \(\Omega_{a,b}=\RR^{a}\times[0,\infty)^{b}\)，\(F\) 为 \(d\)-维紧致纤维），且 \(\pi\) 为投影到基底坐标。
将 \(\omega\) 写作“基方向 \(\wedge\) 纤维方向”的双分次形式，并对纤维变量做 Fubini 积分；对纤维方向应用带角的 Stokes 定理得到
\[
d(\pi_!\omega)-(-1)^d\pi_!(d\omega)=\text{（纤维边界贡献的推前）},
\]
即为 \((\pi^{\partial})_!(\omega|_{\partial_vE})\)。
由于 \(\omega\in\Omega_{\mathrm{rel}}^{k}(E)\) 在 \(\partial_hE\) 上为零，水平边界不产生贡献；角情形下垂直边界分解为各面之和，其选择性消失对应“仅单侧面贡献”的机制（命题 \ref{prop:cdim-stokes-half-face-count}）。
\par
对 \(\rho\) 的专门化：垂直边界由各塌缩面 \(r_j=0\) 组成，其同调类由每个塌缩圆周方向的相对 Thom 类生成；
而这些 Thom 类在定理 \ref{thm:cdim-stokes-explicit-relative-cohomology-basis} 中由 \([\Omega_{\{j\}}]\) 给出，故缺陷项落入其张成子空间。证毕。
\end{proof}
\leanverified{paper\_cdim\_stokes\_fiber\_integration\_defect}

\begin{theorem}[复合纤维积分的二阶 Stokes 缺陷分解]\label{thm:cdim-stokes-composite-fiber-integration-defect}
设 \(\pi:E\to B\)、\(\sigma:B\to C\) 为满足本文假设的带角流形上的定向、适当光滑子沉，
纤维维数分别为 \(d_\pi,d_\sigma\)。
对相对形式 \(\omega\in\Omega^{d_\pi+d_\sigma+k}_{\mathrm{rel}}(E)\)，在 \(\Omega^{k+1}(C)\) 中有精确恒等式
\[
\boxed{
\begin{aligned}
d\big(\sigma_!(\pi_!\omega)\big)
&=(-1)^{d_\pi+d_\sigma}\,(\sigma\circ\pi)_!(d\omega)
\;+\;(-1)^{d_\sigma}\,\sigma_!\big((\pi^{\partial})_!(\omega|_{\partial_vE})\big)
\;+\;(\sigma^{\partial})_!\big((\pi_!\omega)|_{\partial_vB}\big).
\end{aligned}}
\]
其中 \(\sigma^{\partial}:\partial_vB\to C\) 为 \(\sigma\) 的垂直边界诱导映射，两个边界项分别对应于 \(E\) 与 \(B\) 的垂直边界贡献，并在该层级上完全分离。
\end{theorem}
\begin{proof}
对 \(\pi\) 应用命题 \ref{prop:cdim-stokes-fiber-integration-defect} 得
\[
d(\pi_!\omega)=(-1)^{d_\pi}\,\pi_!(d\omega)+(\pi^{\partial})_!(\omega|_{\partial_vE}).
\]
再对 \(\sigma\) 应用同一命题并取 \(\eta:=\pi_!\omega\in\Omega^{d_\sigma+k}_{\mathrm{rel}}(B)\)，得到
\[
d(\sigma_!\eta)=(-1)^{d_\sigma}\,\sigma_!(d\eta)+(\sigma^{\partial})_!(\eta|_{\partial_vB}).
\]
将前式代入并利用无边界主项的复合兼容性 \(\sigma_!(\pi_!(d\omega))=(\sigma\circ\pi)_!(d\omega)\)，整理符号即得所示恒等式。证毕。
\end{proof}

\endinput


\begin{proposition}[径向塌缩纤维化与 Stokes 转移同构]\label{prop:cdim-stokes-radial-fibration-transgression}
\leanverified{paper\_cdim\_stokes\_radial\_fibration\_transgression}
在定理 \ref{thm:cdim-signed-stokes-cohomological-characterization} 的识别
\(
\mathfrak X_0(M)\cong \TT^{u}\times \overline{\mathbb D}^{\,v}
\)
下，写 \(\overline{\mathbb D}^{\,v}\ni z=(z_1,\dots,z_v)\) 并定义径向映射
\[
\rho:\mathfrak X_0(M)\longrightarrow \TT^{u}\times[0,1]^{v},
\qquad
(\theta,z)\longmapsto \bigl(\theta,|z_1|,\dots,|z_v|\bigr).
\]
则 \(\rho\) 在开稠密层 \(\TT^{u}\times(0,1]^{v}\) 上为局部平凡纤维丛，典型纤维为 \(\TT^{v}\)；
当且仅当某一径向坐标 \(|z_j|=0\) 时，对应角变量在纤维上塌缩，故 \(\rho\) 以立方体面格为分层给出一族 torus-塌缩纤维化。
并且在该分层纤维化诱导的谱序列中，
纤维上同调 \(H^{1}(\TT^{v};\ZZ)\) 的全部生成元经转移同构到相对群
\(
H^{2}(\mathfrak X_0(M),\Sigma_0(M);\ZZ)\cong \ZZ^{v}
\)；
在低阶上，该转移同构与命题 \ref{prop:cdim-stokes-exact-sequence-dictionary} 的连接同态 \(\delta\) 一致。
\end{proposition}
\begin{proof}
当 \(z_j\neq 0\) 时极坐标给出局部坐标 \((r_j,\phi_j)\in(0,1]\times\TT\)，从而 \(\rho\) 局部同胚于
\(
\TT^{u}\times(0,1]^{v}\times \TT^{v}\to \TT^{u}\times(0,1]^{v}
\)，得局部平凡性。
当某 \(r_j=0\) 时角变量 \(\phi_j\) 消失且纤维维数下降，分层由立方体面格控制。
\par
谱序列的低阶转移与配对 \((\mathfrak X_0(M),\Sigma_0(M))\) 的长正合列在同一过滤上兼容；
而 \(H^{1}(\mathfrak X_0(M))\) 不含纤维 \(\TT^{v}\) 的贡献（定理 \ref{thm:cdim-signed-stokes-cohomological-characterization}），
故纤维 \(H^{1}\) 只能经转移进入相对 \(H^{2}\) 并由秩比较得到同构，且与 \(\delta\) 的坐标化描述一致。证毕。
\end{proof}

\begin{proposition}[torus-塌缩纤维的 Stokes 转移与圆维的相对 \texorpdfstring{$\dd$}{d}Rham 基]\label{prop:cdim-torus-collapse-relative-forms}
沿用命题 \ref{prop:cdim-stokes-radial-fibration-transgression} 的记号。
则存在一族相对闭的 \((v{+}1)\)-形式
\[
\Theta_1,\dots,\Theta_v\in \Omega^{v+1}\bigl(\mathfrak X_0(M),\Sigma_0(M)\bigr),
\qquad \dd\Theta_j=0,
\]
使得对纤维圆周的标准基 \([\gamma_1],\dots,[\gamma_v]\in H_1(\TT^v;\ZZ)\) 及其经连接同态
\(
\delta:H_1(\TT^v;\ZZ)\to H_2(\mathfrak X_0(M),\Sigma_0(M);\ZZ)
\)
的像，成立 Kronecker 配对
\[
\boxed{\;\bigl\langle[\Theta_i],\,\delta[\gamma_j]\bigr\rangle=\delta_{ij}.\;}
\]
因此 \(\{[\Theta_j]\}_{j=1}^v\) 形成 \(\delta(H_1(\TT^v))\subset H_2(\mathfrak X_0(M),\Sigma_0(M);\ZZ)\) 的对偶基，并以纯 Stokes 配对方式实现“圆维”方向的可积核验。
\end{proposition}
\begin{proof}
在开稠密层 \(\TT^{u}\times(0,1]^v\) 上取极坐标 \((r_j,\phi_j)\in(0,1]\times\TT\)，其中 \(\phi_j\) 为纤维 \(\TT^v\) 的第 \(j\) 个角坐标。
取 \([0,1]\) 上紧支撑 \(1\)-形式 \(\alpha_j\)（支撑包含于 \((0,1)\)）并归一化为 \(\int_{0}^{1}\alpha_j=1\)。
令
\[
\Theta_j:=\dd r_j\wedge \dd\phi_j\wedge \bigwedge_{i\neq j}\alpha_i(r_i),
\]
并将其沿 \(\TT^u\) 方向作平凡拉回；由构造其在塌缩边界 \(\Sigma_0(M)\) 的邻域相对消失，故定义为相对形式。
显然 \(\dd\Theta_j=0\)。
\par
连接同态 \(\delta[\gamma_j]\) 可由“沿 \(r_j\) 扫出一张 \(2\)-膜并在 \(r_j=0\) 处塌缩”的相对链实现；在局部模型中该链同胚于 \([0,1]_{r_j}\times S^1_{\phi_j}\)。
于是
\[
\int_{\delta[\gamma_j]}\Theta_i
=\begin{cases}
\int_{0}^{1}\!\!\int_{0}^{2\pi}\dd r_j\,\dd\phi_j\cdot \prod_{k\ne j}\int_0^1\alpha_k=1,& i=j,\\[4pt]
0,& i\ne j,
\end{cases}
\]
即得到所示 Kronecker 配对。证毕。
\end{proof}
\leanverified{paper\_cdim\_torus\_collapse\_relative\_forms}

\begin{conjecture}[缺陷系数的 Stokes 通量刻画]\label{conj:cdim-defect-entropy-stokes-flux}
设 \(\pi\) 为满足本文“扩展保存/外置预算”纪律的可审计投影系统。
存在一个自然伴随的有限生成可消去交换幺半群 \(M_\pi\)，使得其 Stokes 角色配对
\((\mathfrak X_0(M_\pi),\Sigma_0(M_\pi))\)
为 \(\pi\) 的边界/体积接口提供一条同调化记账口径，并且与 \(\pi\) 相容的对数尺度缺陷系数
\(\kappa_\pi\) 满足
\[
\boxed{\;\kappa_\pi=\rank_{\ZZ}H^{2}\bigl(\mathfrak X_0(M_\pi),\Sigma_0(M_\pi);\ZZ\bigr)\;}
\]
从而缺陷系数在该口径下由相对上同调的通量格秩量子化。
在正交幺半群扇区 \(M_\pi\cong \ZZ^{u}\oplus \NN^{v}\oplus F\) 中，上式退化为 \(\kappa_\pi=v\)（定理 \ref{thm:cdim-signed-stokes-cohomological-characterization}）。
\end{conjecture}

\paragraph{审计实现的半整数预算下界：几何嵌入不可降低结构分账}
为把“实现层可编码性”与“可审计可逆解码”在同一结构计量中区分，考虑如下实现纪律。

\begin{definition}[审计实现与圆维预算]\label{def:cdim-audited-realization-budget}
令 \(M\) 为满足定义 \ref{def:cdim-signed} 的幺半群。
称注入幺半群同态
\[
\iota:\ M\hookrightarrow \RR^{a}\times (\RR_{\ge 0})^{b}\times F'
\]
为一个\textbf{审计实现}，若满足：
\begin{enumerate}
  \item \(\iota\) 在群完成层诱导注入
        \(
        G(M)\hookrightarrow \RR^{a+b}\times F'
        \)
        且像为离散子群；
  \item \(\iota(U(M))\subseteq \RR^{a}\times\{0\}^{b}\times F'\)（单位元方向不得落入单侧坐标）。
\end{enumerate}
定义该实现的\textbf{圆维预算}为
\[
\boxed{\;B(\iota):=a+\frac{b}{2}\;}
\]
（半整数取值）。
\end{definition}

\begin{theorem}[审计实现的最小圆维预算下界]\label{thm:cdim-audited-realization-budget-lower-bound}
在定义 \ref{def:cdim-audited-realization-budget} 的设定下，任意审计实现均满足
\[
\boxed{\;B(\iota)\ \ge\ \cdim^{\pm}(M)\;}
\]
并且对 \(M=\ZZ^{u}\oplus \NN^{v}\oplus F\) 的正交情形，下界可达。
\end{theorem}
\begin{proof}
由 \(\iota(U(M))\subseteq \RR^{a}\times\{0\}^{b}\times F'\) 与离散性，\(U(M)\) 的无挠秩不超过 \(a\)，从而
\(
\cdim(U(M))\le a
\)（例 \ref{ex:circle-dimension-examples}）。
同理，将群完成注入投影到 \(\RR^{b}\) 坐标并模去 \(U(M)\) 的像，可得 \(G(M)/U(M)\) 的无挠秩不超过 \(b\)，故
\(
\cdim(G(M)/U(M))\le b
\)。
代入定义 \ref{def:cdim-signed} 得
\(
\cdim^{\pm}(M)\le a+\frac12 b=B(\iota)
\)。
\par
当 \(M=\ZZ^{u}\oplus \NN^{v}\oplus F\) 时，取标准嵌入
\(
\iota_0:\ZZ^{u}\oplus \NN^{v}\oplus F\hookrightarrow \RR^{u}\times (\RR_{\ge 0})^{v}\times F
\)
即可达到 \(B(\iota_0)=u+v/2=\cdim^{\pm}(M)\)（命题 \ref{prop:cdim-signed-orthant-closed}）。证毕。
\end{proof}
\leanverified{paper\_cdim\_audited\_realization\_budget\_lower\_bound}

\paragraph{相位压缩的精度负担：低维环面注入的分离度指数律}\label{para:cdim-phase-precision-tradeoff}
在把多维离散寄存器压入有限个相位圆时，结构维数缺口并不会消失，而会以分离度的多项式衰减速率进入精度账本。

\begin{definition}[分离度与精度负担指数]\label{def:cdim-phase-separation-precision-burden}
令 \(\TT^{d}:=(\RR/\ZZ)^{d}\)，并取其 \(L^\infty\) 商度量
\[
d_\infty(x,y):=\min_{m\in\ZZ^{d}}\|x-y-m\|_\infty.
\]
对任意映射 \(f:\ZZ^{r}\to \TT^{d}\) 与窗口半径 \(N\ge 1\)，定义分离度
\[
\mathrm{sep}_{f}(N)
:=\min\bigl\{d_\infty(f(x),f(y)):\ x\neq y,\ \|x\|_\infty\le N,\ \|y\|_\infty\le N\bigr\}.
\]
定义其\textbf{精度负担指数}为
\[
\mathrm{pb}(f):=\liminf_{N\to\infty}\frac{-\log \mathrm{sep}_{f}(N)}{\log N}\in[0,\infty].
\]
\end{definition}

\begin{theorem}[相位--精度权衡不等式]\label{thm:cdim-phase-precision-tradeoff}
对任意注入 \(f:\ZZ^{r}\hookrightarrow \TT^{d}\)，有
\[
\boxed{\;\mathrm{pb}(f)\ \ge\ \frac{r}{d}\;}
\]
从而存在常数 \(C>0\) 使对无穷多 \(N\) 成立
\(
\mathrm{sep}_{f}(N)\le C\,N^{-r/d}
\)。
\leanverified{paper\_cdim\_phase\_precision\_tradeoff}
\end{theorem}
\begin{proof}
记 \(B_N:=\{x\in\ZZ^{r}:\ \|x\|_\infty\le N\}\)，则 \(\card{B_N}=(2N+1)^{r}\)。
将 \(\TT^{d}\) 以 \(L^\infty\) 边长 \(1/M\) 的网格立方体覆盖，所需立方体数至多为 \(M^{d}\)。
取 \(M:=\left\lfloor(2N+1)^{r/d}\right\rfloor\)。
当 \(N\) 足够大时有 \(M^{d}<\card{B_N}\)，故由鸽巢原理存在 \(x\neq y\in B_N\) 使 \(f(x),f(y)\) 落入同一网格立方体，从而
\(
d_\infty(f(x),f(y))\le 1/M\le 2(2N+1)^{-r/d}
\)。
因此 \(\mathrm{sep}_{f}(N)\le 2(2N+1)^{-r/d}\) 对无穷多 \(N\) 成立，代入定义 \ref{def:cdim-phase-separation-precision-burden} 即得 \(\mathrm{pb}(f)\ge r/d\)。证毕。
\end{proof}

\begin{corollary}[dyadic 深度的必需增长率]\label{cor:cdim-phase-precision-tradeoff-dyadic}
在定理 \ref{thm:cdim-phase-precision-tradeoff} 的设定下，若要求在窗口 \(B_N\) 上由 dyadic 深度 \(b\) 的相位量化（网格尺度 \(2^{-b}\)）实现可分离读出，则必须满足
\[
2^{-b}\ <\ \mathrm{sep}_{f}(N)
\quad\Longrightarrow\quad
b\ \ge\ \frac{r}{d}\log_2 N - O(1)
\qquad(N\to\infty).
\]
\end{corollary}
\begin{proof}
由 \(2^{-b}<\mathrm{sep}_{f}(N)\) 与定理 \ref{thm:cdim-phase-precision-tradeoff} 的 \(\mathrm{sep}_{f}(N)\le C N^{-r/d}\)（无穷多 \(N\)）合并，取对数即得。证毕。
\end{proof}
\leanverified{paper\_cdim\_phase\_precision\_tradeoff\_dyadic}

\begin{remark}[实现层单射与审计预算的区分]\label{rem:cdim-coding-vs-audit}
命题 \ref{prop:impl-coding} 表明：只要对象与证书具有有限描述，即存在可计算单射把其实现层序列化为自然数。
该事实仅陈述“实现层可存储”，并不产生可用于改变结构圆维的自由度。
在审计纪律下（离散晶格/可分离读出），结构圆维缺口必以预算或精度指数形式进入外置账本：
定理 \ref{thm:cdim-audited-realization-budget-lower-bound} 给出几何实现预算的半整数下界，
定理 \ref{thm:cdim-phase-precision-tradeoff} 给出相位压缩的分离度指数下界。
这种“结构圆因子/残差轴”的分离可由圆维签名 \(\Sigma\Sigma\) 统一记录（定义 \ref{def:cdim-signature}），并在可逆素集合约束下由 \(S\)-solenoid 终对象性质给出实现无关的最小因子化（定理 \ref{thm:cdim-s-solenoid-terminal-object}）。
\end{remark}

\subsubsection{直积自然数半环的同态刚性与幂等分割通道}\label{subsubsec:cdim-semiring-rigidity}
在定义 \ref{def:cdim-audited-realization-budget} 的审计实现口径中，单侧正锥常以 \(\NN^v\) 的坐标半轴作为“互斥寄存器通道”的抽象。
当协议不仅要求保加法，还要求兼容乘法组合规则时，自然出现交换幺半环同态的语义限制。
下面给出一条刚性结论：在 \(\NN^r\) 的坐标逐点乘加结构下，同态不允许任何混合编码，其可做之事仅为坐标抽取（允许重复）。
该刚性为“实现层可编码”与“结构层不可坍缩”提供一个完全代数化的分界线。

\begin{theorem}[直积自然数半环的同态刚性]\label{thm:cdim-nr-nd-semiring-hom-rigidity}
将 \(\NN^r,\NN^d\) 视为交换幺半环：加法与乘法均按坐标逐点定义，且 \(0=(0,\dots,0)\)、\(1=(1,\dots,1)\)。
则任意交换幺半环同态
\[
f:\NN^r\longrightarrow \NN^d
\]
都存在唯一函数 \(\sigma:\{1,\dots,d\}\to\{1,\dots,r\}\) 使得对一切 \(x=(x_1,\dots,x_r)\in\NN^r\) 有
\[
f(x)=\bigl(x_{\sigma(1)},x_{\sigma(2)},\dots,x_{\sigma(d)}\bigr).
\]
\leanverified{semiring\_hom\_rigidity}
\leanverified{paper\_cdim\_nr\_nd\_semiring\_hom\_rigidity}
\end{theorem}
\begin{proof}
记 \(\pi_j:\NN^d\to\NN\) 为第 \(j\) 个坐标投影，并令 \(f_j:=\pi_j\circ f:\NN^r\to\NN\)。
则每个 \(f_j\) 仍是交换幺半环同态。
令
\[
e_i:=(0,\dots,0,1,0,\dots,0)\in\NN^r
\]
为第 \(i\) 个坐标单位向量，则 \(e_i^2=e_i\)，且当 \(i\neq k\) 时 \(e_ie_k=0\)，并且 \(\sum_{i=1}^r e_i=1\)。
\par
在 \(\NN\) 中满足 \(n^2=n\) 的仅有 \(n\in\{0,1\}\)，故 \(f_j(e_i)\in\{0,1\}\)。
由同态性，
\[
0=f_j(0)=f_j(e_ie_k)=f_j(e_i)\,f_j(e_k),
\]
于是对固定 \(j\)，至多存在一个 \(i\) 使 \(f_j(e_i)=1\)。
另一方面，
\[
1=f_j(1)=f_j\!\Bigl(\sum_{i=1}^r e_i\Bigr)=\sum_{i=1}^r f_j(e_i),
\]
故必恰有一个指标 \(i=\sigma(j)\) 使 \(f_j(e_{\sigma(j)})=1\)，其余为 \(0\)。
对任意 \(x=\sum_{i=1}^r x_ie_i\)（按加法分解）有
\[
f_j(x)=\sum_{i=1}^r x_i f_j(e_i)=x_{\sigma(j)}.
\]
汇总各 \(j\) 得所述形式。
唯一性由每个 \(j\) 的指标 \(\sigma(j)\) 唯一确定。证毕。
\leanverified{paper\_circleDim\_axiomatic\_completeness}
\end{proof}

\begin{corollary}[单射与自同构的结构限制]\label{cor:cdim-nr-nd-semiring-hom-inj-auto}
在定理 \ref{thm:cdim-nr-nd-semiring-hom-rigidity} 的设定下：
\begin{enumerate}
  \item 若 \(f:\NN^r\to\NN^d\) 为单射交换幺半环同态，则 \(\sigma\) 必为满射，特别地 \(d\ge r\)；
  \item 若 \(f\) 为同构，则 \(d=r\) 且 \(\sigma\) 为置换，从而
  \(
  \Aut_{\mathrm{Semiring}}(\NN^r)\cong S_r
  \)。
\end{enumerate}
\leanverified{semiring\_aut\_is\_perm}
\end{corollary}
\begin{proof}
(1) 若存在 \(i_0\notin\sigma(\{1,\dots,d\})\)，则改变输入的第 \(i_0\) 坐标不改变输出，故 \(f\) 非单射。
反之若 \(\sigma\) 满射，则任意 \(x\neq y\) 至少在某一输入坐标 \(i\) 上不同；取 \(j\) 使 \(\sigma(j)=i\)，则输出第 \(j\) 坐标不同，故 \(f\) 单射。
(2) 同构必单射且满射；由定理 \ref{thm:cdim-nr-nd-semiring-hom-rigidity} 的显式形式立即得到 \(d=r\) 且 \(\sigma\) 双射，从而为置换。证毕。
\end{proof}

\begin{proposition}[幂等分割刻画 \(\Hom(\NN^r,S)\)]\label{prop:cdim-nr-hom-idempotent-splitting}
令 \(S\) 为任意交换幺半环。
则交换幺半环同态集合 \(\Hom_{\mathrm{Semiring}}(\NN^r,S)\) 与集合
\[
\Bigl\{(\varepsilon_1,\dots,\varepsilon_r)\in S^r:\ 
\varepsilon_i^2=\varepsilon_i,\ 
\varepsilon_i\varepsilon_k=0\ (i\neq k),\
\sum_{i=1}^r \varepsilon_i=1
\Bigr\}
\]
之间存在自然双射，其对应由
\[
f\ \longmapsto\ (f(e_1),\dots,f(e_r))
\]
给出；逆向由
\[
(\varepsilon_1,\dots,\varepsilon_r)\ \longmapsto\
f_{\varepsilon}(x_1,\dots,x_r):=\sum_{i=1}^r x_i\varepsilon_i
\]
给出。
\leanverified{semiring\_hom\_determines\_idempotents}
\leanverified{paper\_cdim\_nr\_hom\_idempotent\_splitting}
\end{proposition}
\begin{proof}
若 \(f\) 为同态，则 \(f(e_i)^2=f(e_i^2)=f(e_i)\)，且 \(f(e_ie_k)=f(0)=0\) 给出正交性；
又 \(f(\sum e_i)=f(1)=1\) 给出和为 \(1\)。
\par
反之，给定正交幂等分割 \((\varepsilon_i)\) 并定义 \(f_{\varepsilon}\) 如上，则加法同态性显然且 \(f_\varepsilon(1)=\sum \varepsilon_i=1\)。
乘法同态性由正交性保证：
\[
f_\varepsilon(x)\,f_\varepsilon(y)
=\sum_{i,k} x_iy_k\,\varepsilon_i\varepsilon_k
=\sum_i x_iy_i\,\varepsilon_i
=f_\varepsilon(x\cdot y),
\]
其中 \(x\cdot y\) 表示逐点乘积。
两映射互为逆由构造直接核验。证毕。
\end{proof}

\begin{conjecture}[乘法通道强制的最小半圆维预算]\label{conj:cdim-semiring-idempotent-budget}
令 \(S\) 为交换幺半环，并设 \(S\) 含有 \(r\) 个两两正交且和为 \(1\) 的非零幂等元 \((\varepsilon_1,\dots,\varepsilon_r)\)。
则在定义 \ref{def:cdim-audited-realization-budget} 的“可审计实现”纪律的乘法兼容版本中，任何实现都必须提供至少 \(r\) 个互相独立的可审计账本道，且其带符号圆维预算满足
\[
B(\iota)\ \ge\ \frac{r}{2},
\]
并且等号当且仅当该实现与把 \(r\) 个幂等通道实现为 \(r\) 个互斥单侧寄存器的标准正交模型等价。
\end{conjecture}

\endinput


\endinput
